\section{Why the scale is \texorpdfstring{$t^\kappa$}{t to the kappa}}\label{S:why-kappa}

This section explains, without proof, where the exponent $\kappa$ and the form
of $C_2(\Gamma,\alpha_0)$ in Theorem~\ref{T:intro-second-moment} come from.

\begin{figure}[htbp]
  \centering
  % Exponent only; open markers denote excluded parameter endpoints.
\begin{tikzpicture}[x=4.7cm,y=2.7cm,font=\normalsize]
  \draw[ham/arrow] (0,1)--(1.06,1);
  \draw[ham/arrow] (0,1)--(0,2.04) node[above] {$\kappa$};
  \foreach \xx/\lab in {0/0,.5/{1/2},1/1} {
  \draw[hamblue] (\xx,1)--(\xx,.96);
  \node[anchor=north] at (\xx,.94) {$\lab$};
  }
  \foreach \yy/\lab in {1/1,{3/2}/{3/2},{5/3}/{5/3},2/2} {
  \draw[hamblue] (-.025,{\yy})--(0,{\yy});
  \node[anchor=east] at (-.04,{\yy}) {$\lab$};
  }
  \draw[hamblue] (-.025,{7/5})--(0,{7/5});
  \draw[hamblue,line width=.45pt]
  (-.025,{7/5})--(-.055,{7/5})--(-.055,1.30);
  \node[anchor=east,inner sep=0pt] at (-.07,1.30) {$7/5$};
  \node[align=center,text width=4.1cm] at (.5,1.93)
  {Static limit at left ends:\\ $(4-\alpha)/(3-\alpha)$};
  \draw[hamblue,line width=1pt] (0,{7/5})--(1,1);
  \draw[hamteal,line width=1pt] (0,{3/2})--(1,1);
  \draw[hamgold,line width=1pt] (0,{5/3})--(.5,{4/3});
  \draw[hamgold!45,densely dashed,line width=.8pt]
  (.5,{4/3})--(1,1);
  \foreach \yy/\colour in {{7/5}/hamblue,{3/2}/hamteal,{5/3}/hamgold} {
  \filldraw[draw=\colour,fill=white,line width=.8pt] (0,{\yy}) circle (1.5pt);
  }
  \draw[hamgold!65,dotted,line width=.45pt]
  (.5,{4/3})--(.56,1.40);
  \node[anchor=south west,inner sep=0pt,text=hamgold] at (.56,1.40) {$4/3$};
  \filldraw[draw=hamgold,fill=white,line width=.8pt] (.5,{4/3}) circle (1.5pt);
  \filldraw[draw=black!65,fill=white,line width=.8pt] (1,1) circle (1.5pt);
  \draw[hamblue,line width=1pt] (1.12,1.96)--(1.27,1.96);
  \node[anchor=west] at (1.32,1.96) {$\alpha=1/2$};
  \draw[hamteal,line width=1pt] (1.12,1.75)--(1.27,1.75);
  \node[anchor=west] at (1.32,1.75) {$\alpha=1$};
  \node[anchor=west] at (1.12,1.56) {white in space};
  \node[anchor=west] at (1.12,1.39) {$d=1$};
  \draw[hamgold,line width=1pt] (1.12,1.17)--(1.27,1.17);
  \node[anchor=west] at (1.32,1.17) {$\alpha=3/2$};
  \node[anchor=west,align=left] at (1.12,.87) {dashed: outside\\ hypothesis};
  \node at (.5,.61) {temporal exponent $\alpha_0$};
  \draw[ham/arrow] (.27,.40)--(.10,.40);
  \node[anchor=west,inner sep=0pt] at (.31,.40) {longer memory};
\end{tikzpicture}

  \caption[The exponent $\kappa$ against the temporal exponent]{The exponent
    $\kappa=1+(1-\alpha_0)/(3-\alpha)$ of~\eqref{E:kappa}, drawn against the
    temporal exponent $\alpha_0$ for three spatial exponents $\alpha$. Longer
    memory in time (smaller $\alpha_0$) and a more singular spatial covariance
    (larger $\alpha$) both raise $\kappa$. As $\alpha_0\downarrow0$ each line
    tends to the static exponent $(4-\alpha)/(3-\alpha)$. The end points are
    excluded: Hypothesis~\ref{H:homog} requires $0<\alpha_0<1$ and
    $\alpha+\alpha_0<2$, so the line for $\alpha=3/2$ stops at $\alpha_0=1/2$.
    The picture shows the exponent only; it says nothing about the constant
    $C_2(\Gamma,\alpha_0)$.}
  \label{Fig:kappa-fan}
\end{figure}


The first step is an exact scaling, proved in
Lemma~\ref{L:exact-canonical-scaling}. Stretch time by a factor $T$ and space
by $r_T=T^{(\alpha_0-1)/(3-\alpha)}$. Because the temporal covariance
$|s-t|^{-\alpha_0}$ and the spatial covariance $\Gamma$ are homogeneous
(Hypothesis~\ref{H:homog}), the stretched noise has the law of the original
noise multiplied by a constant. With $\delta=T^{-\kappa}$, the stretched
solution has the law of the solution $V_\delta$ of
\[
  \delta^2\partial_t^2V_\delta=\Delta V_\delta+\sqrt\delta\,\W V_\delta,
  \qquad V_\delta(0,\cdot)=1,\quad \partial_tV_\delta(0,\cdot)=0,
\]
and
\[
  T^{-\kappa}\log\E|u(T,0)|^2=\delta\log\E|V_\delta(1,0)|^2
\]
(see~\eqref{E:exact-fast-equation} and~\eqref{E:exact-moment-scaling}). With
this choice of $\kappa$, the time derivative carries the factor $\delta^2$ and
the noise the factor $\sqrt\delta$. The growth of the second moment over long
times thus becomes a question about small noise: does
$\delta\log\E|V_\delta(1,0)|^2$ have a limit as $\delta\to0$?

Heuristically it does, because a gain and a cost of the same order $1/\delta$
balance. Shift the noise toward a deterministic potential, chosen through a
nonnegative control $\Psi$, so that the potential seen by the wave at time $s$
is $h_{\Psi,s}$ of~\eqref{E:intro-control-mean}. In the scaled equation this
needs a shift of order $1/\sqrt\delta$, and its Gaussian cost is $v_\Psi/(2\delta)$,
one half of the variance~\eqref{E:intro-control-variance} divided by $\delta$.
Freezing the potential at time $s$ gives the equation
$\delta^2a''=\mathcal{R}_\Psi(s)\,a$ for the amplitude of the leading mode,
whose growing solution has exponential rate $\sqrt{\mathcal{R}_\Psi(s)}/\delta$.
Heuristically, the square of the wave then grows like
$\exp\bigl(2\int_I\sqrt{\mathcal{R}_\Psi(s)}\ud s/\delta\bigr)$. Gain minus
cost, multiplied by $\delta$, is
$2\int_I\sqrt{\mathcal{R}_\Psi(s)}\ud s-v_\Psi/2$, and taking the supremum over $\Psi\in\mathcal{C}_{\mathrm{pos}}$ gives
$C_2(\Gamma,\alpha_0)$. Returning to the original variables,
$\log\E u^2(T,x)$ grows like $C_2(\Gamma,\alpha_0)\,T^\kappa$.

This argument only explains the answer. A proof needs a lower bound for each
control, which computes the cost using the original covariance, retaining correlations between different time intervals
(Chapter~\ref{S:book-lower}), and an upper bound showing that no other
fluctuation of the noise does better (Chapters~\ref{S:book-means}--\ref{S:book-point-upper});
Chapter~\ref{S:book-constant} matches the two. Section~\ref{S:strategy} describes
these steps.
